\section{Quality Evaluation and Quantization}


\begin{questionbox}{Domanda 49}
\textbf{Compare MSE/PSNR, SSIM and LPIPS}: what each measures, pros/cons, and why two images with the same MSE can have very dif and only iferent perceived quality.
\end{questionbox}


\begin{itemize}
  \item \textbf{MSE} $\rightarrow$ average squared pixel error:

\end{itemize}
$$
\text{MSE}=\frac{1}{NM}\sum_{i,j}(x_{ij}-\hat{x}_{ij})^2
$$

\begin{itemize}
  \item \textbf{PSNR} $\rightarrow$ logarithmic MSE:

\end{itemize}
$$
\text{PSNR}=10\log_{10}\left(\frac{MAX^2}{\text{MSE}}\right)
$$

\begin{itemize}
  \item \textbf{MSE/PSNR} $\rightarrow$ simple and reproducible, but weak perceptually; ignore structure, masking, semantics.
  \item \textbf{SSIM} $\rightarrow$ compares local luminance, contrast, structure.
  \item \textbf{LPIPS} $\rightarrow$ compares deep features; better for perceptual/semantic distortions, but expensive and model-dependent.
  \item \textbf{Same MSE can look dif and only iferent} $\rightarrow$ error distribution matters: texture noise less visible; edge/structure error more annoying.

\end{itemize}
\begin{questionbox}{Domanda 50}
(MMQuEv) Explain the dif and only iference between subjective and objective quality evaluation in multimedia systems and why both are necessary.
\end{questionbox}


\begin{itemize}
  \item \textbf{Subjective evaluation} $\rightarrow$ human observers rate perceived quality.
  \begin{itemize}
    \item Ground truth for perception.
    \item Slow, expensive, statistically variable.
  \end{itemize}
  \item \textbf{Objective evaluation} $\rightarrow$ automatic signal/feature metric.
  \begin{itemize}
    \item Fast, repeatable, scalable.
    \item May correlate imperfectly with human perception.
  \end{itemize}
  \item \textbf{Need both} $\rightarrow$ subjective validates perception; objective enables optimization and monitoring.

\end{itemize}
\begin{questionbox}{Domanda 51}
(MMQuEv) What are the key stages involved in designing a subjective quality test according to standardized guidelines?
\end{questionbox}


\begin{itemize}
  \item \textbf{Define}:
  \begin{itemize}
    \item dataset/content selection;
    \item display/audio equipment;
    \item viewing/listening conditions;
    \item \textbf{methodology}: ACR, DSIS, pairwise comparison, etc.;
    \item participant screening;
    \item score processing/outlier handling.
  \end{itemize}
  \item \textbf{Control factors} $\rightarrow$ viewing distance, brightness/contrast, room illumination, audio setup, test duration.

\end{itemize}
\begin{questionbox}{Domanda 52}
(MMQuEv) Describe the main categories of objective quality metrics based on the availability of the original “reference” signal.

(MMQuEv) Which of the following best describes the “Full-Reference” (FR) objective quality assessment approach?
\end{questionbox}


\begin{itemize}
  \item \textbf{Full Reference (FR)} $\rightarrow$ original and degraded signals available; examples: MSE, PSNR, SSIM, VMAF.
  \item \textbf{Reduced Reference (RR)} $\rightarrow$ only compact original features available.
  \item \textbf{No Reference (NR)} $\rightarrow$ only degraded signal available; examples: BRISQUE, NIQE, PIQE.
  \item \textbf{FR luminance PSNR}:

\end{itemize}
$$
\text{MSE}_Y(t)=\frac{1}{NM}\sum_{n,m}[I(n,m,1,t)-\hat{I}(n,m,1,t)]^2
$$

$$
\text{PSNR}_Y(t)=10\log_{10}\left(\frac{255^2}{\text{MSE}_Y(t)}\right)
$$

\begin{questionbox}{Domanda 53}
(MMQuEv) In the context of subjective testing, what is the main purpose of a “screening” phase for participants?
\end{questionbox}


\begin{itemize}
  \item \textbf{Screening} $\rightarrow$ checks participants can correctly perceive stimuli.
  \item \textbf{Removes invalid observers}: uncorrected visual issues, color-vision problems, inconsistent scoring behavior.

\end{itemize}
\begin{questionbox}{Domanda 54}
(MMQuEv) Why is statistical analysis a critical component of subjective quality evaluation?
\end{questionbox}


\begin{itemize}
  \item \textbf{Human scores vary} $\rightarrow$ need reliability measures.
  \item \textbf{Statistical analysis} $\rightarrow$ MOS, variance, confidence intervals, outlier detection.
  \item \textbf{MOS}:

\end{itemize}
$$
\text{MOS}=\frac{1}{N}\sum_{i=1}^{N}x_i
$$

\begin{itemize}
  \item \textbf{Standard error}:

\end{itemize}
$$
SE=\frac{s}{\sqrt{N}}
$$

\begin{questionbox}{Domanda 55}
(S\&P Qnt) Explain the dif and only iference between a “mid-tread” and a “mid-rise” quantizer in the context of uniform quantization for signed data.
\end{questionbox}


\begin{itemize}
  \item \textbf{Mid-tread} $\rightarrow$ zero is reconstruction level; small values around zero map to zero.
  \item \textbf{Mid-rise} $\rightarrow$ zero is decision threshold; values around zero map to positive/negative non-zero levels.
  \item \textbf{Residuals prefer mid-tread} $\rightarrow$ many small values become exactly zero:

\end{itemize}
$$
Q(x)=\Delta\cdot\text{round}\left(\frac{x}{\Delta}\right)
$$

\begin{questionbox}{Domanda 56}
(S\&P Qnt) Define the concept of a “deadzone” in a quantizer and explain why it is frequently employed in lossy compression systems.
\end{questionbox}


\begin{itemize}
  \item \textbf{Deadzone quantizer} $\rightarrow$ enlarged interval mapped to zero:

\end{itemize}
$$
Q(x)=0 \quad \text{for small } |x|
$$

\begin{itemize}
  \item \textbf{Useful in lossy compression} $\rightarrow$ prediction/transform residuals often have many small coefficients.
  \item \textbf{Consequence} $\rightarrow$ more zeros $\rightarrow$ better entropy coding.

\end{itemize}
\begin{questionbox}{Domanda 57}
(S\&P Qnt) Why is scalar quantization alone often considered insufficient for effective compression of non-sparse data?
\end{questionbox}


\begin{itemize}
  \item \textbf{Scalar quantization} $\rightarrow$ one sample at a time.
  \item \textbf{Non-sparse signal} $\rightarrow$ many significant samples remain $\rightarrow$ entropy coding has little to exploit.
  \item \textbf{Better compression needs}:
  \begin{itemize}
    \item \textbf{prediction} $\rightarrow$ reduce residual variance;
    \item \textbf{transform coding} $\rightarrow$ compact energy;
    \item \textbf{vector/block coding} $\rightarrow$ exploit dependencies;
    \item \textbf{entropy coding} $\rightarrow$ exploit non-uniform probabilities.

  \end{itemize}
\end{itemize}
\begin{questionbox}{Domanda 58}
(S\&P Qnt) What is the condition for a predictive quantization system to be effective, and how is the “coding gain” defined?
\end{questionbox}


\begin{itemize}
  \item \textbf{Effective condition}:

\end{itemize}
$$
\sigma_y^2 < \sigma_x^2
$$

\textbf{with}:

$$
y(n)=x(n)-v(n)
$$

\begin{itemize}
  \item \textbf{Prediction gain}:

\end{itemize}
$$
G_P=10\log_{10}\left(\frac{\sigma_x^2}{\sigma_y^2}\right)
$$

\begin{itemize}
  \item Positive when predictor reduces variance.

\end{itemize}
\begin{questionbox}{Domanda 59}
(S\&P Qnt) Draw the scheme of a linear predictive quantization scheme, and motivate the structure, with particular attention to the use of a decoding loop at the encoder side.
\end{questionbox}


\begin{itemize}
  \item Predictor must use reconstructed past samples on encoder and decoder:

\end{itemize}
$$
v(n)=\mathcal{P}(\hat{x}(n-1),\hat{x}(n-2),\dots)
$$

$$
y(n)=x(n)-v(n)
$$

$$
\hat{x}(n)=\hat{y}(n)+v(n)
$$

\begin{itemize}
  \item Encoder includes same reconstruction loop as decoder $\rightarrow$ prevents drift.

\end{itemize}
\begin{questionbox}{Domanda 60}
(S\&P Qnt) What is the primary purpose of the predictor in a predictive quantization system?

(S\&P Qnt) In a predictive quantization system, if the prediction $v(n)$ is nearly equal to $x(n)$, what happens to the variance of the signal $y(n)$ being sent to the quantizer?
\end{questionbox}


\begin{itemize}
  \item \textbf{Predictor} $\rightarrow$ exploits correlation among neighboring samples.
  \item If $v(n)\approx x(n)$:

\end{itemize}
$$
y(n)=x(n)-v(n)\approx 0
$$

$$
\sigma_y^2 \ll \sigma_x^2
$$

\begin{itemize}
  \item Quantizer encodes lower-energy residual.

\end{itemize}
\begin{questionbox}{Domanda 61}
(S\&P Qnt) When selecting a linear predictor of order $P$ for a random process, how does the prediction error variance typically behave as the order increases?
\end{questionbox}


\begin{itemize}
  \item \textbf{Linear predictor}:

\end{itemize}
$$
v(n)=-\sum_{i=1}^{P}a_i x(n-i)
$$

$$
y(n)=\sum_{i=0}^{P}a_i x(n-i),\quad a_0=1
$$

\begin{itemize}
  \item Increasing $P$ cannot worsen optimal error variance in theory $\rightarrow$ old solution still available.
  \item \textbf{Practice} $\rightarrow$ larger $P$ increases complexity, side information, estimation error; gains eventually negligible.

\end{itemize}
\begin{questionbox}{Domanda 62}
\textbf{Explain the screening effect}: why does the prediction gain saturate as the linear predictor order $P$ increases?
\end{questionbox}


\begin{itemize}
  \item Increasing $P$ $\rightarrow$ optimal prediction error variance cannot increase in theory.
  \item Gain usually saturates due to \textbf{screening effect}.
  \item Nearby samples already explain most correlation with current sample.
  \item Farther samples add little independent information because correlation is mostly mediated by nearer samples.
  \item \textbf{Result} $\rightarrow$ diminishing returns plus higher complexity/estimation sensitivity.

\end{itemize}
\begin{questionbox}{Domanda 63}
(S\&P Qnt) In high-resolution uniform quantization, what is the approximate relationship between the SNR and the bit rate ($R$)?
\end{questionbox}


\begin{itemize}
  \item High-resolution uniform quantization:

\end{itemize}
$$
D \approx \frac{A^2}{12}2^{-2R}
$$

$$
\text{SNR}
=10\log_{10}\left(\frac{\sigma_X^2}{D}\right)
=10\log_{10}\left(\frac{\sigma_X^2}{\frac{A^2}{12}2^{-2R}}\right)
$$

\begin{itemize}
  \item With $\gamma^2=\frac{A^2}{4\sigma_X^2}$:

\end{itemize}
$$
\text{SNR}\approx 6.02R - 10\log_{10}\left(\frac{\gamma^2}{3}\right)
$$

\begin{itemize}
  \item \textbf{Rule} $\rightarrow$ each extra bit/sample gives about $6$ dB SNR improvement.

\end{itemize}
---

